% Section 1 — Introduction
% Drafted by the lead from docs/method.md, docs/AMENDMENT-A1.md and the verified reference base.
% No result values appear here: outcome-dependent sentences live in Section 7.

\section{Introduction}
\label{sec:intro}

Event-driven architectures have become the default way to move money between financial
institutions. A payment is ingested as an event, validated, split into debit and credit legs, and
appended to a general ledger; downstream services notify counterparties, update balances, and
trigger settlement. Apache Kafka is the most common substrate for this pipeline, and its
\emph{exactly-once semantics} (EOS) is routinely presented as the guarantee that makes the
pipeline safe: a record is said to be processed exactly once even in the presence of broker
failures, consumer crashes and rebalances~\cite{kip98,kafka-design}.

The guarantee is real, but bounded. Kafka's transactions make \emph{Kafka-internal} effects
atomic: records appended to Kafka topics, consumer offsets committed through the transaction
coordinator, and the producer's own idempotency state all commit or abort together. When the
useful side effect of the computation must land in an external system --- a PostgreSQL ledger, a
core-banking database, a payment rail --- that write is a separate transaction at a separate site,
with its own write-ahead log and no shared coordinator. Kafka's own documentation is explicit that
in this case the consumer's position must be coordinated with what is actually stored as
output~\cite{kafka-design}; the exactly-once guarantee is not a distributed transaction spanning
Kafka and the external store~\cite{kip98,txncommit2006}.

For most applications the resulting duplicates or gaps are an annoyance to be cleaned up later. In
a double-entry ledger they are a correctness failure of a different kind. Accounting represents
every economic event as a journal entry whose debit and credit legs sum to zero, and it requires
that the ledger be append-only and that each business event appear exactly once. A duplicate
posting is not a redundant row that can be deleted in a reconciliation job without breaking
auditability: it silently changes balances, often in a way that still sums to zero and therefore
passes naive balance checks. A lost posting is worse: the money exists on one side of the books
only, and no later retry can be distinguished from a legitimate new event unless the pipeline
carried an idempotency mechanism. Both failures are silent, permanent, and exactly the kind of
defect that surfaced in industrial incident reports and in the correctness discussions around
streaming data pipelines~\cite{kafka2015,flink2017,bailis2014}.

Practitioners nevertheless deploy pipelines in which Kafka transactions are treated as if they
covered the database write. The pattern has a long history: an application consumes a record,
writes to the database, and then commits the offset inside a Kafka transaction (or, in older
Spring code, inside a ``chained'' transaction manager), with the implicit assumption that either
both commits happen or neither does. When the process dies between the two commits, or when one
commit succeeds while the other times out, the assumption fails. What is missing from the
literature is not the theoretical observation that two sites cannot be committed atomically
without a protocol such as two-phase commit~\cite{txncommit2006,pg-prepare}, but a systematic,
empirical mapping from failure modes to ledger-invariant violations, mechanism by mechanism, with
the cost of each remedy measured under the same conditions. Survey work on microservice data
consistency documents the patterns and the practitioner confusion but reports no fault-injection
counts~\cite{laigner2021,micro2026}; systems papers on exactly-once semantics stop at the
framework's own state~\cite{flink2017,millwheel2013,dstreams2013}; and ledger/auditability work
assumes the invariant rather than testing its preservation~\cite{triple2019,triple2021}.

This paper closes that gap for the canonical case: a Kafka source topic whose records are
balanced journal entries, a consumer service that must append the corresponding postings to an
external relational database and maintain consumption progress, and a pipeline that must keep
four ledger invariants under scripted faults. We ask three questions.
\textbf{RQ1:} for which failure modes does Kafka EOS fail to preserve double-entry ledger
invariants when the sink is an external relational database?
\textbf{RQ2:} how do five mechanisms and three combinations compare --- naive at-least-once
(M0, negative control), Kafka EOS with a Kafka-internal sink (M1-kafka, positive control), Kafka
EOS wrapped around a non-atomic external write (M1x, the deployed pattern), a transactional
outbox/inbox (M2), idempotency keys with a unique index (M3), offsets stored in the ledger
database transaction (M4), and M2+M3, M4+M3, M1x+M3 --- on duplicate postings, lost postings,
invariant violations, throughput and tail latency?
\textbf{RQ3:} what is the minimal mechanism set sufficient for each failure-exposure class?

Our approach has four components. First, we define the ledger state machine and the precise
\emph{EOS boundary} of the pipeline, enumerate every point at which a process can crash or a
message can be lost, duplicated or reordered (crash points C1--C12), and give an
invariant-preservation proposition or an explicit counterexample trace for every mechanism and
crash-point class. A bounded version of the crash-point semantics is machine-checked with TLC,
which confirms the expected duplicate and loss witnesses within stated bounds; the model is a
bounded witness, not an unbounded proof, and we are explicit about the difference. Second, we
implement all nine configurations in one code base with identical tuning, and inject six fault
modes with real scripted injectors --- broker kills, consumer kills, produce timeouts with unknown
outcome, commit-order failures in both directions, database network partitions, and duplicate
redelivery after rebalance --- alongside deterministic crash-point failpoints. Third, we score
every run with an independent oracle: the load generator writes the expected posting set before
producing, and a separate Python checker recomputes it from the specification and compares it with
the database after quiescence, so duplicates, losses, imbalance and mutation are measured rather
than inferred from logs. Fourth, we pre-register the statistical analysis plan (exact confidence
intervals for mostly-zero correctness outcomes, bootstrap intervals and Friedman/Nemenyi/Cliff's
delta for performance) and execute a staged campaign designed to fit a fixed wall-clock budget,
recording every deviation.

Our contributions are:
\begin{itemize}[leftmargin=1.1em,itemsep=1pt]
\item a formal ledger and EOS-boundary model with a crash-point enumeration and
      invariant-preservation propositions for nine mechanism configurations, including a
      machine-checked bounded TLC model of the duplicate (W1) and loss (W2) windows;
\item the first fault-injection comparison, to our knowledge, of EOS-boundary mechanisms scored
      directly on double-entry invariants with an independent oracle and pre-registered statistics;
\item a practitioner decision table mapping failure exposure to the minimal sufficient mechanism
      set, with the measured cost of each additional guarantee;
\item a public, reproducible artifact containing the service, the fault injectors, the oracle,
      the raw run records and the scripts that regenerate every table and figure.
\end{itemize}

The remainder of the paper is organised as follows. Section~\ref{sec:related} reviews Kafka EOS,
consistency mechanisms, ledger invariants and experimental methodology. Section~\ref{sec:model}
defines the system, ledger state machine, invariants and the EOS boundary.
Section~\ref{sec:mechanisms} presents the mechanisms. Section~\ref{sec:formal} states and argues
the invariant-preservation results and reports the machine-checked model.
Section~\ref{sec:method} describes the implementation, faults, oracle and pre-registered design.
Section~\ref{sec:results} reports correctness and performance results and the hypothesis verdicts.
Section~\ref{sec:decision} derives the decision table, Section~\ref{sec:threats} discusses threats
to validity, and Section~\ref{sec:conclusion} concludes.
